\chapter{Đúng và Sai}
\markboth{Đúng và Sai}{Right and Wrong}

\section{Đúng lý nhưng sai cách nói.}
\begin{truyen}{Đúng lý nhưng sai cách nói.}{Being right is not enough if love is absent.}
\chuhoa{A}{nh biết em mình sai nên nói rất thẳng. Cái sai thì gọi đúng tên rồi, nhưng cái giọng trên cơ làm người nghe chỉ muốn bật lại. Nhiều khi mình đúng nội dung mà hỏng sạch ở cách nói, thành ra điều đúng cũng mất tác dụng.}
\textit{(He knew his younger sibling was wrong and spoke bluntly. The content was correct, but the tone was full of contempt, so the listener only wanted to resist. Some things are right in content yet wrong in delivery, and then even truth cannot save the relationship.)}
\end{truyen}
\begin{baihoc}
Đúng mà nói kiểu dằn mặt thì cũng thành hỏng việc.
\textit{(Rightness without humanity can become another kind of wrong.)}
\end{baihoc}
\begin{ghinhoanh}
\textbf{Short English to remember:}

Truth also needs tenderness.
\end{ghinhoanh}
\begin{tuhocnhanh}
1. Em có lần nào nói đúng nhưng làm người khác đau theo cách không cần thiết không? \textit{(Have you ever spoken the truth in a way that hurt more than necessary?)}

2. Em có thể giữ điều đúng mà đổi cách nói như thế nào? \textit{(How can you keep the truth while changing the way you speak it?)}
\end{tuhocnhanh}
\ngancach

\section{Sai một lần nhưng biết sửa ngay.}
\begin{truyen}{Sai một lần nhưng biết sửa ngay.}{A repaired mistake is better than hidden pride.}
\chuhoa{E}{m chép nhầm dữ liệu trong bài nhóm, làm cả nhóm ngồi sửa mệt nghỉ. Thay vì đổ tại file, tại người này người kia, em nhận luôn, thức tối sửa lại và xin lỗi từng bạn. Sai thì vẫn là sai, nhưng cách gánh nó lên mới nói ra mình là người kiểu gì.}
\textit{(The student copied data incorrectly and caused the whole group to lose time. Instead of blaming others, the mistake was admitted, repaired that night, and each teammate was personally apologized to. The mistake remained a mistake, but the way it was faced determined growth.)}
\end{truyen}
\begin{baihoc}
Sai mà chịu sửa còn hơn đúng vài lần rồi sinh ra ngông.
\textit{(Being wrong and correcting it can be more valuable than being right and becoming proud.)}
\end{baihoc}
\begin{ghinhoanh}
\textbf{Short English to remember:}

Repair matters more than pride.
\end{ghinhoanh}
\begin{tuhocnhanh}
1. Khi sai, em thường che hay sửa? \textit{(When wrong, do you usually hide or repair?)}

2. Sai lầm gần đây nhất của em đang cần một hành động cụ thể nào? \textit{(What concrete action does your latest mistake require?)}
\end{tuhocnhanh}
\ngancach

\section{Đúng theo số đông nhưng sai với lương tâm.}
\begin{truyen}{Đúng theo số đông nhưng sai với lương tâm.}{Crowds can approve what conscience rejects.}
\chuhoa{C}{ả lớp xúm vào chọc một bạn nên nhìn bề ngoài cứ như chuyện bình thường. Em cũng cười hùa cho đỡ lạc quẻ, nhưng trong bụng biết rõ là bẩn. Đông người đứng cùng một bên không có nghĩa bên đó tử tế.}
\textit{(The whole class mocked one student, so the behavior felt normal. The student laughed along just to avoid standing apart, yet inside knew it was not right. Some things approved by the crowd are simply many people being wrong together.)}
\end{truyen}
\begin{baihoc}
Đám đông đông tới đâu cũng không thay được lương tâm.
\textit{(The measure of the crowd cannot replace the measure of conscience.)}
\end{baihoc}
\begin{ghinhoanh}
\textbf{Short English to remember:}

The crowd can be wrong together.
\end{ghinhoanh}
\begin{tuhocnhanh}
1. Em có từng làm điều mình biết không đúng chỉ vì ai cũng làm? \textit{(Have you ever done something you knew was wrong simply because everyone did it?)}

2. Lương tâm em đang nhắc điều gì mà em chưa nghe đủ? \textit{(What is your conscience telling you that you have not listened to enough?)}
\end{tuhocnhanh}
\ngancach

\section{Sai vì thiếu hiểu biết khác với sai vì cố ý.}
\begin{truyen}{Sai vì thiếu hiểu biết khác với sai vì cố ý.}{Not every wrong comes from the same place.}
\chuhoa{M}{ột học sinh nói sai vì thật sự chưa hiểu. Một học sinh khác biết rõ mà vẫn quay cóp. Cả hai đều sai, nhưng sai kiểu khác nhau thì không thể quăng chung một rọ. Muốn xử cho công bằng thì phải nhìn đúng gốc của cái sai, chứ không chỉ nhìn bề mặt.}
\textit{(One student answered wrongly because the lesson was not yet understood. Another knew the truth but still cheated. Both were wrong, yet they could not be treated the same because the roots were different. Justice lies not only in naming the wrong, but in understanding its nature correctly.)}
\end{truyen}
\begin{baihoc}
Muốn trị đúng bệnh thì đừng chẩn đoán kiểu qua loa.
\textit{(To repair well, we must correctly see the root of the wrong.)}
\end{baihoc}
\begin{ghinhoanh}
\textbf{Short English to remember:}

Different wrongs need different responses.
\end{ghinhoanh}
\begin{tuhocnhanh}
1. Khi thấy người khác sai, em có vội kết luận quá nhanh không? \textit{(When seeing someone wrong, do you judge too quickly?)}

2. Em cần phân biệt rõ hơn điều gì giữa lỗi và ý xấu? \textit{(What do you need to distinguish more clearly between mistake and bad intent?)}
\end{tuhocnhanh}
\ngancach

\section{Sống đúng không phải để hơn người khác.}
\begin{truyen}{Sống đúng không phải để hơn người khác.}{Right living is not a stage for ego.}
\chuhoa{C}{ó người làm đúng vài chuyện rồi nhìn ai cũng thấy thấp hơn mình. Họ không biết cái đúng của mình đang bị dùng để nuôi cái tôi. Sống đúng mà càng lúc càng khinh người thì coi chừng cái đúng đó đã lệch ngay từ trong ruột.}
\textit{(Some people do what is right and then look at others with superiority. They do not realize that their rightness is being used as capital to feed the ego. When right living produces pride, that rightness has already begun to tilt.)}
\end{truyen}
\begin{baihoc}
Đúng mà làm mình bớt kiêu mới là đúng tới nơi tới chốn.
\textit{(What is right is most beautiful when it makes us humbler, not more arrogant.)}
\end{baihoc}
\begin{ghinhoanh}
\textbf{Short English to remember:}

Right living should soften pride.
\end{ghinhoanh}
\begin{tuhocnhanh}
1. Em có đang dùng điều đúng của mình để phán xét người khác không? \textit{(Are you using your rightness to judge others?)}

2. Làm sao để sống đúng mà vẫn giữ được lòng khiêm tốn? \textit{(How can you live rightly while remaining humble?)}
\end{tuhocnhanh}
\ngancach
